\begin{answers}{ch:policy-gradient}{方策勾配法：REINFORCE から PPO まで}
\item 分散の合計は、$b=0$ で $\frac{4}{27}\approx0.1481$、$b=\frac13$ で $\frac{2}{27}\approx0.0741$ です。
$\vtheta=\vzero$ なので $\bm{\pi}=(\frac13,\frac13,\frac13)$ です。
$b=0$ では、報酬があるのは $a=1$ だけなので $\E\norm{\hat{\vg}}^2=\frac13\times1^2\times\frac69=\frac29$ です。
平均は厳密な勾配 $\E\hat{\vg}=\frac13(\ve_1-\bm{\pi})=(\frac29,-\frac19,-\frac19)$ で、$\norm{\E\hat{\vg}}^2=\frac{6}{81}=\frac{2}{27}$ です。
よって分散は $\frac29-\frac{2}{27}=\frac{4}{27}$ です。
$b=\frac13$ では $(r(a)-b)^2$ が $\frac49,\frac19,\frac19$ なので、$\E\norm{\hat{\vg}}^2=\frac13\cdot\frac69\cdot\frac69=\frac{4}{27}$ です。
平均はベースラインで変わらず $\frac{2}{27}$ のままなので、分散は $\frac{4}{27}-\frac{2}{27}=\frac{2}{27}$ です。
実測の 0.1487 と 0.0741 は、20 万サンプルの誤差の範囲で一致します。
\item (a) $\E_{a\sim\pi_{\text{old}}}[\rho]=\sum_a\pi_{\text{old}}(a\mid s)\frac{\pitheta(a\mid s)}{\pi_{\text{old}}(a\mid s)}=\sum_a\pitheta(a\mid s)=1$ です。
(b) 期待値は足し算に分けられるので、$\E[\rho-1-\log\rho]=\E[\rho]-1-\E[\log\rho]=-\E[\log\rho]$ です（(a) を使いました）。
$-\log\rho=\log\frac{\pi_{\text{old}}}{\pitheta}$ なので、$-\E_{\pi_{\text{old}}}[\log\rho]$ は KL ダイバージェンスの定義そのもの $\KL{\pi_{\text{old}}}{\pitheta}$ です。
(c) $f(x)=x-1-\log x$ とおくと $f'(x)=1-\frac1x$ です。
$f'$ は $0<x<1$ で負、$x>1$ で正なので、$f$ は $x=1$ で最小になり、最小値は $f(1)=0$ です。
したがって 1 サンプルごとの値が 0 以上になり、推定が負になることがありません。
\item $\hat{A}_0=0.388$、$\hat{A}_1=0.36$、$\hat{A}_2=0.2$ です。
$\gamma$ が変わると TD 誤差 $\delta_t$ そのものが変わるので、\tabref{tab:pg-gae-example}の $\delta_t$ をそのまま使ってはいけません（使うと $0.264$、$0.28$ という誤った値になります）。
$\gamma=1$ で計算し直すと、$\delta_2=1+0-0.8=0.2$、$\delta_1=0+0.8-0.6=0.2$、$\delta_0=0+0.6-0.5=0.1$ です。
漸化式 $\hat{A}_t=\delta_t+\gamma\lambda\hat{A}_{t+1}$（$\gamma\lambda=0.8$）を後ろから当てはめると、
$\hat{A}_2=0.2$、$\hat{A}_1=0.2+0.8\times0.2=0.36$、$\hat{A}_0=0.1+0.8\times0.36=0.388$ です。
\item 推定はすべてのサンプルで $\vzero$ になり、平均も $\vzero$ です（実測。20 万サンプルで平均・分散とも 0）。
厳密な勾配 $(0.222,-0.111,-0.111)$ とは違うので、偏った推定です。
$r(a)-b=r(a)-r(a)=0$ なので、どの行動を引いても重みが 0 になるからです。
\eqref{eq:pg-baseline-zero}の証明では、$b$ が行動によらないので $b$ を内側の期待値の外に出し、$\E_a[\nabla\log\pitheta(a)]=\vzero$ を使いました。
$b=r(a)$ は行動 $a$ で決まるので外に出せず、$\E_a[r(a)\nabla\log\pitheta(a)]$ が残ります。これは厳密な勾配そのもので、0 ではありません。
\item 1 回集めたデータで 1 回だけ勾配を計算する、GAE 付きのアクター・クリティック（A2C の形。ベースライン付き REINFORCE で $G_t-V$ を GAE に変えたもの）と同じになります。
更新は集めた直後の 1 回だけなので、確率比はいつも $\rho=1$ で、クリップは一度も効きません（$\epsilon$ を変えても結果は同じです）。
実測（各 3 個の種）では、最後の 10 反復の平均エピソード長は 275 で、本文の PPO の 494 に遠く及びませんでした。
40 反復（640 エピソード）の時点でも約 120 で、PPO の約 470 より大幅に遅れます。
PPO の速さは、同じデータで 40 回（10 エポック × 4 ミニバッチ）更新することから来ているのです。
$\epsilon$ を変えた結果は、本文の\figref{fig:pg-eps-lambda}(a)(b) にあります（各 3 個の種、40 反復。\file{sweep_ppo_eps.py}）。
最後の 10 反復の平均長（3 個の種の平均）は、$\epsilon=0.05$ で 466、0.2 で 453、0.5 で 389 でした。
1 反復で方策が動く量（KL の平均）は、$\epsilon=0.05$ で約 0.0004、0.5 で約 0.025 です。
$\epsilon$ は大きいほど 1 回の反復で動ける幅が広くなり、0.5 では種によって 339〜468 とばらつき、クリップなし（$\epsilon=\infty$）の不安定さに近づきます。
小さくすると動ける幅は狭くなりますが、この易しい問題では、0.05 でも最後の長さは下がりませんでした。

\item reward-to-go は分散をまったく減らしません。
報酬が最後の 1 回だけで $\gamma=1$ なら、どの時刻でも「その時刻から後の報酬の和」は最後の報酬そのもので、$G_t=R(\tau)$ です（本文の式\eqref{eq:pg-llm}の直後でも、同じことを述べました）。
つまり、reward-to-go の推定は、REINFORCE の推定と 1 サンプルごとに同じ値になります。
reward-to-go が減らしていたのは「その行動より前の報酬」による揺らぎで、ここではそれが最初から 0 だからです。
残る手段はベースラインです。価値ネットを使わない限り、1 つの応答のトークンはすべて同じ重み $R(\tau)-b$ を受け取ります。すべてのトークンが同じ $R(\tau)$ を分け合うので、トークンごとの価値を細かく学ぶ利点は小さくなります。
一方、トークンごとの価値ネットは LLM と同じくらい大きく、学習も難しく、メモリも 2 倍近く要ります。
そこで、同じプロンプトに複数の回答を出し、その報酬の平均をベースラインにする「回答ごとのベースライン」がよく使われます（\chapref{ch:rl-llm}）。
\end{answers}
